\documentclass{article}

\usepackage[T1]{fontenc} 
\usepackage[french]{babel}
\usepackage[utf8]{inputenc}
\usepackage{graphicx}
\graphicspath{{./img/}}
\usepackage{amsmath, amssymb}
\usepackage[left=3cm,right=3cm,top=2cm,bottom=2cm]{geometry}
\usepackage{url}
\usepackage[linktocpage=true]{hyperref}

\title{Rapport de Projet : Tower Defense}
\author{Youri CHANCRIN, Théo DESCOMPS, Samuel KHALIFA, Jean-Baptiste MARTINEZ}
\date{\today}

\makeatletter
\let\mytitle\@title
\let\myauthor\@author
\makeatother

\begin{document}

\begin{flushleft}

\thispagestyle{empty}
\includegraphics[width=0.2\textwidth]{enseirb-matmeca}
\graphicspath{{./img/}}

\vspace{\stretch{0.4}}

\hrulefill \\[2em]
\begin{center}
\textbf{\Huge \mytitle}\\[0.8em]
{\large equipe : projetss6-tower-19048}
\\[0.8em]
\end{center}
\hrulefill

\vspace{\stretch{0.6}}

{\large \textbf{Auteurs :} \myauthor}\\
\vspace{\stretch{0.2}}
{\large \textbf{Responsable de projet :} David Renault}\\
\vspace{\stretch{0.2}}
{\large \textbf{Enseignant référent :} Frédéric Herbreteau}

\vspace{\stretch{1}}

\begin{center}
  Première année, filière informatique\\
  Date : \today
\end{center}

\vspace{\stretch{1}}

\hfill Version actuelle du document : draft

\end{flushleft}

\newpage

\tableofcontents

\newpage

\section{Sujet}
L'objectif de ce projet est la réalisation d'un jeu de type tower défense.

Ce genre peut se résumer en une seule phrase :\\
\center \label{description_jeu} Des attaquants envahissent une zone défendue par des défenseurs.

La programmation fonctionnelle sera mise en avant.
L'executable sera réalisé en Typescript.

\subsection{Theme}
Nous avons choisi de créer un jeu de tower defense autour du thème
du Pastafirsme.
Selon cette religion, la température de la Terre est inversement proportionnelle 
au nombre de pirates qu’elle contient.
Il faut donc répandre cette idéologie en défendant le monstre spaghetti volant et 
en augmentant le nombre de pirate sur terre.

\section{Conception}

Débutons la conception avec notre description du jeu \ref{description_jeu} 
pour en tirer des élements fondamentaux.

On remarque en premier lieu la présence d'entité, des attaquants, mais aussi des 
défenseurs. Bien qu'elles soient antagonistes, ces entités ont des similarités, notamment 
dans le traitement de leurs actions. 

On décide donc de les regrouper sous le même nom, les acteurs.


\subsection{Entité Acteur}

Les acteurs forment un ensemble d'entités qui intéragissent entre elles.
Pour définir ces interactions, chaque acteur a besoin d'avoir un ensemble 
de comportements. Pour que ces interactions prennent effet, nos entités ont 
égalements besoin d'un état interne.

Le type acteur est très vague et peut représenter beaucoup de choses.
Pour garder une consistence entre nos acteurs, on définit des archétypes.

(Bon la partie qui suit je sais pas si c'est pas plus de la réalisation :/, 
peut être réduire ou déplacer ? besoin d'avis)

\begin{itemize}
   \item Spaghetti Monster : Le gentil monstre spaghetti que les défenseurs doivent protéger.
   \item Good Guy : Les défenseurs.
   \item Ignorant : Un type d'attaquants.
   \item Ignorance Spreader : Un autre type d'attaquants.
   \item Spawner : Un point d'apparition pour les attaquants.
   \item Ground : Un bout de chemin sur lequel les attaquans se déplacent.
\end{itemize}

Les types Good Guy, Ignorant et Ignorance Spreader représentent 
les défenseurs et les attaquants, ils découlent de notre démarche.
Cependant, nos acteurs comportent aussi un type Spawner et Ground.
Ces types ne peuvent pas être attaquer ou attaquer, mais ils ont 
d'autres interactions avec le reste des entités qui justifient qu'il 
soient aussi considérés comme des acteurs.

En effet, Spawner peut être vu comme un attaquant immortel, et Ground 
un défenseur qui contraint le déplacement des attaquants.

Avec la diversification des acteurs vient un nouveau problème,
les attributs spécifiques.

Dans le cas de la \textbf{position}, par exemple, tous les acteurs en ont besoin.
Mais si l'on s'intéresse aux attributs spécifiques aux combats, les acteurs 
qui n'attaquent pas n'ont pas d'intérêts à les avoir.

On pourrait rajouter tous les attributs nécessaires à tous les comportements 
à tous les acteurs. Mais c'est une mauvaise idée pour deux raisons.
Premièrement, cela entrainerait une modification sur le type acteur.
Ce-dernier étant utilisé dans de nombreux contextes, cela favoriserait 
l'apparition de comportements imprévus.
Secondement, le type acteur aurait toujours plus d'états possibles, 
et deviendrait de plus en plus complexe et lourd à comprendre et gérer.

Pour éviter ces désagréments, nous avons séparer la responsabilités des attributs 
dépendants de certains comportements dans un champs \textbf{externalProps}.
Pour que ce champs puisse satisfaire toutes les possibilités d'attributs, il
est de type dictionnaire de quelconque. Toutefois, ce typage faible est compensé 
par la présence de fonctions getter et setter agissants comme une interface 
avec \textbf{externalProps}.

Nos acteurs sont donc parés à toutes éventalité d'évolution.
Cependant, bien qu'ils disposent de comportements, ils ne peuvent pas se modifier 
entre eux directement.
Nous avons besoin de quelque chose d'ordre supérieur pour gérer leurs interactions.
Quelque chose comme un système de phase.

\subsection{Phase}




\section{Boucle de jeu}
- Un monde = plan 2d
- acteurs = séparé du monde.
Phases:
 - Acteur propose une action
 - Phase transforme les acteurs en fonction de leurs propositions.
 - Puis, le monde impose ses limites, corrigeant les acteurs invalides (en dehors du monde)


\section{Phases, typage fort ou typage faible}
Choix : Typage fort
Parce que : Impose rigourosité et cohérence entre les types des propositions des 
acteurs et la fonction transformant les acteurs et permet d'expérimenter avec 
le riche système de type de Typescript, notamment les mapped types

Inconvénient : pas dynamique, mais pas un besoin pour nous donc pas important

- Actions par défaut.

\section{Structure du monde et des acteurs}
Choix : Les acteurs sont séparés du monde
A cause de l'habitude de la programmation object, on aurait envie d'établir le plus de liens possible entre nos structures
Parce qu'on est dans la programmation fonctionnelle, on ne peut modifier l'attribut d'un 'object' dans le recréer, 
donc avoir les acteurs dans le monde est très impratique, car les acteurs sont sans cesse recréé, ceux qui forcerait de 
recréer le monde pour rien.

\section{Position des acteurs}
Chaque acteur contient sa position.
Si un acteur ne bouge pas, il n'est pas modifié.

Inconvénient : complexité n carré, n le nombre d'acteur, pour avoir tous les acteurs à une certaine distance d'un acteur.

Alternative : Faire un tableau 2D avec chaque acteur dans une case. Cependant, ce n'est pas une structure fonctionnelle, 
le projet portant sur cette dernière, on évite donc l'utilisation d'une matrice des positions.

\section{Génération aléatoire des Acteurs}
- Génération aléatoire des chemins
- Génération aléatoire des Acteurs

\section{Détail de la structure Actor}
- externalProps

\section{Règles et fin de partie}

\section{limite pureté}
- Dès qu'on fait du random, c'est plus pure
- Probleme de perf ?

\end{document}